\documentclass[12pt,a4paper,oneside]{article}
\usepackage[brazil]{babel}
\usepackage[utf8]{inputenc}
\usepackage{olymp}

\setcounter{problem}{3}

\begin{document}

\begin{problem}{Metade da prova}{prova}{2}

José é um excelente aluno e acerta todas as questões que faz. Mas ele também é preguiçoso e não quer fazer a prova toda, já que só precisa de metade dos pontos para passar. Ele faz as questões na ordem e para quando já completa exatamente a metade dos pontos da prova. Por exemplo, suponha uma prova com 4 questões, de valores 5, 3, 2 e 10, nesta ordem. Após terminar a questão 3 ele já tem 10 pontos ($5+3+2$), que é a metade do valor total de prova, então ele para de fazer a prova na questão 3.

Sua tarefa é, dada os valores das questões, descobrir até qual questão da prova José fará.


%\Notes
%
%Observações, se tiver.

\InputFile

A entrada é composta por duas linhas. A primeira contém um valor inteiro $N$, que indica o número de questões da prova. A segunda linha contém $N$ valores inteiros $q_1, q_2, \dots, q_N$ indicando o valor de cada questão. Restrições: $1 < N \leq 10^5$; $1 \leq q_i \leq 100, \forall i=1\dots N$.

\OutputFile

Seu programa deve imprimir uma única linha contendo um inteiro que indica a número da última questão que José vai fazer.

\Notes
É garantido que sempre existe uma resposta.

\Example

\begin{example}
\exmp{
4
5 3 2 10
}{
3
}%
\end{example}

\begin{example}
\exmp{
9
2 8 2 8 4 4 4 4 4
}{
4
}%
\end{example}

% \small \textit{Obs.: baseada na questão ``Guerra por território'' da OBI 2012} 

%Comentários sobre o exemplo, se necessário.

\end{problem}

\end{document}
